\chapter{Implementation Details}
\label{ch:progress}

This chapter documents the implementation status and key engineering decisions made during development.

\section{Overall Implementation Status}
\label{sec:impl_status}

\noindent\textbf{Live Deployment:} The system is currently deployed at \url{http://techpulse-dev-frontend-939514668437.s3-website-us-east-1.amazonaws.com/}.

Table~\ref{tab:impl_status} summarizes the implementation status of each system component as of the reporting date.

\noindent\textbf{Status key:} \textit{Unit-tested} = implemented and verified by unit tests only; \textit{Integration-tested} = verified end-to-end in local Docker Compose environment; \textit{Cloud-validated} = deployed and verified in the AWS Free Tier.

\begin{table}[htbp]
\caption{Component Implementation Status Summary.}
\label{tab:impl_status}
\centering
\begin{tabular}{p{4.2cm}p{2.8cm}p{5.5cm}}
\toprule
\textbf{Component} & \textbf{Status} & \textbf{Notes} \\
\midrule
Data Ingestion (5 sources) & Integration-tested & ArXiv, Hacker News, DEV.to, GitHub, RSS fully implemented \\
Preprocessing Pipeline & Integration-tested & 17-step normalization + token-based chunking \\
Embedding Pipeline & Integration-tested & MiniLM L6 v2 (384-dim), lazy singleton \\
Hybrid Retrieval & Cloud-validated & 5-stage pipeline (metadata filter, vector, BM25, weighted RRF, cross-encoder) \\
Baseline Retrieval & Integration-tested & Vector-only comparison mode \\
RAG Orchestrator & Integration-tested & 3-layer hallucination verification \\
LLM Backend Fallback & Unit-tested & Auto-cascade across 4 backends \\
API Layer & Integration-tested & FastAPI + Mangum; /health, /ask, /drift endpoints \\
Database Schema & Integration-tested & PostgreSQL + pgvector with HNSW index \\
S3 Medallion Storage & Integration-tested & Raw/processed/embeddings tiers (LocalStack locally) \\
SQS Queue + DLQ & Integration-tested & Ingestion--embedding decoupling (LocalStack locally) \\
Drift Detection & Unit-tested & Dual-criteria: 10\% threshold + Shewhart $3\sigma$ \\
Health Monitoring & Integration-tested & 4-probe deep health check (DB, S3, SQS, LLM) \\
CloudWatch Integration & Cloud-validated & 8 metrics in 1 namespace (\texttt{TechPulse/\{stage\}}) \\
Evaluation Framework & Cloud-validated & Paired 50-query run (baseline + hybrid configurations), RAGAS, statistical tests, and cost projection; follow-on drift-impact analysis remains future work \\
CI/CD Pipeline & Cloud-validated & 4-stage GitHub Actions + SAM \\
SAM Infrastructure & Cloud-validated & Full SAM template; stack deployed and operational \\
Docker Compose (Local) & Integration-tested & 6 services (db, pgadmin, api, frontend, scheduler, localstack) \\
React Frontend & Integration-tested & Query interface with citation display \\
Unit Tests & Integration-tested & 208 tests across 15 test files, all passing \\
\bottomrule
\end{tabular}
\end{table}

\section{Key Enhancements Beyond the Original Proposal}
\label{sec:enhancements}

Two significant enhancements were made beyond the original proposal: (1) data source coverage was expanded from two to five sources (ArXiv, Hacker News, DEV.to, GitHub Trending, RSS feeds) to improve corpus diversity and temporal coverage; and (2) the preprocessing pipeline was extended from basic HTML stripping to a 17-step normalization pipeline (validated by 18 dedicated unit tests) to handle the diverse content formats encountered across sources, particularly Markdown-heavy content from DEV.to and GitHub. All five ingesters share the same idempotent, SHA-256-based deduplication architecture.

Additionally, nine research-backed retrieval improvements were applied post-proposal, including: weighted Reciprocal Rank Fusion (vector~0.50, BM25~0.35, recency~0.15), improved BM25 tokenization with stop-word removal, BM25-only candidate quality gate, asymmetric candidate limits, post-fusion score gate, Shewhart $3\sigma$ minimum standard deviation floor for drift detection, RAGAS evaluation threading fix, cross-encoder reranking (ms-marco-MiniLM-L-6-v2), and NaN retry pass for RAGAS scoring. These improvements are documented in the project's engineering change log and collectively produced a composite quality swing of $+0.069$---from trailing the baseline by $-0.022$ to leading by $+0.047$.

\section{Preprocessing Pipeline Architecture}
\label{sec:preprocessing_detail}

The 17-step text normalization pipeline is applied to all ingested documents before chunking to handle the heterogeneous content formats across five data sources. The steps are sequenced to prevent cascading errors (e.g., stripping HTML tags before removing entity references):

\begin{enumerate}[leftmargin=0.5in, font=\small]
  \item HTML entity unescaping (\texttt{\&nbsp;} $\to$ space)
  \item HTML comment removal (\texttt{<!-- -->})
  \item HTML tag stripping (e.g., \texttt{<p>}, \texttt{<div>}, \texttt{<span>})
  \item Fenced code block removal (\texttt{```...```})
  \item Inline code removal (\texttt{`...`})
  \item Markdown image/link reference removal (\texttt{![alt](url)}, \texttt{[text](url)})
  \item URL removal (http, https, ftp schemes)
  \item Markdown heading removal (\texttt{\# \#\# \#\#\#})
  \item Markdown bold/italic marker removal (\texttt{**}, \texttt{*}, \texttt{\_\_}, \texttt{\_})
  \item Markdown strikethrough removal (\texttt{~~text~~})
  \item Markdown horizontal rule removal (\texttt{---}, \texttt{***})
  \item Markdown blockquote marker removal (\texttt{>})
  \item Emoji and pictographic symbol removal (Unicode ranges for emoticons, transport, supplemental symbols, regional indicators, dingbats, variation selectors)
  \item Unicode control character removal (bytes \texttt{0x00--0x1F}, zero-width joiners, byte-order marks)
  \item Repeated punctuation collapse (e.g., \texttt{!!!} $\to$ \texttt{!})
  \item Unicode whitespace normalization (\texttt{0xA0} $\to$ space)
  \item General whitespace collapse (multiple consecutive spaces, tabs, newlines)
\end{enumerate}

Following normalization, text is split into overlapping token-based chunks using the GPT-4 tokenizer (\texttt{cl100k\_base} encoding from TikToken) with configurable chunk size and overlap. Default parameters are chunk size 500 tokens ($\approx$400 words) and 50-token overlap, balancing context coherence against retrieval candidate diversity. The chunking operation produces 6{,}220 indexed chunks from the 3{,}984 ingested documents.

\section{AWS Free Tier Deployment}
\label{sec:freetier}

TechPulse is deployed on a personal AWS Free Tier account. The deployment adapts certain services from the original proposal to remain within Free Tier cost constraints while preserving full architectural capabilities. 

\subsection*{Cost Breakdown (100 queries/day estimate)}

Table~\ref{tab:cost_breakdown} provides a component-level cost estimate for the anticipated 100 queries/day operational scenario:

\begin{table}[H]
\centering
\caption{Estimated Monthly Cost Breakdown at 100 Queries/Day Workload.}
\label{tab:cost_breakdown}
\begin{tabular}{p{3.2cm}p{2.0cm}p{3.5cm}p{2.0cm}}
\toprule
\textbf{Component} & \textbf{Cost/Month} & \textbf{Details} & \textbf{Free Tier} \\
\midrule
Lambda (ingestion) & \$0.05 & 4 invocations/day (every 6\,h) & \checkmark \\
Lambda (query handler) & \$0.08 & 100 queries/day, 3.1s avg & \checkmark \\
Groq LLM API & \$0.12\textsuperscript{$\dagger$} & $\sim$2,300 tokens/query average & Free (500k/day) \\
RDS PostgreSQL & \$0.08 & db.t3.micro instance (micro is free) & \checkmark \\
S3 + EventBridge + SQS & \$0.00 & Storage, routing, queue services & \checkmark \\
\midrule
\textbf{Total Monthly} & \textbf{\$0.33} & \textbf{Extrapolated from Phase B eval} & \textbf{\checkmark} \\
\bottomrule
\end{tabular}
\end{table}

\noindent\textsuperscript{$\dagger$}Groq inference is free within the 500k tokens/day allowance; \$0.12 represents the imputed cost at equivalent commercial per-token rates to enable fair cost projection if the free tier were unavailable.

\noindent Costs are derived from observed token usage and Lambda invocation metrics collected during the evaluation run. At higher workloads exceeding the Groq free-tier allowance (500,000 tokens/day), the automatic LLM backend fallback chain (Groq $\to$ Bedrock; each backend has one designated fallback partner) enables cost-controlled scaling. The monthly total of \$0.33 is well below the \$15 budget target, providing a 45$\times$ safety margin.

\begin{figure}[htbp]
\centering
\includegraphics[width=0.88\textwidth]{costbreakdownoverview.png}
\caption{AWS cost breakdown overview: monthly costs by component at 100 queries/day workload, demonstrating Free Tier eligibility (\$0.33 projected monthly cost vs.\ \$15 budget).}
\label{fig:cost_breakdown_visual}
\end{figure}

\begin{figure}[htbp]
\centering
\includegraphics[width=\textwidth]{cost_projections.png}
\caption{Cost scaling projections: daily and monthly costs as query volume increases from 50 to 2{,}000 queries/day. Left panel shows the daily variable operating cost (query-proportional Lambda and imputed Groq LLM inference). Right panel shows total monthly cost including fixed infrastructure (Lambda ingestion, RDS db.t3.micro), ranging from \$0.23/month (50~q/day) to \$4.13/month (2{,}000~q/day)---all well within the \$15 budget target. Beyond $\sim$213~queries/day, the Groq free-tier token allowance (500k tokens/day) is exhausted; the automatic LLM fallback chain absorbs overflow.}
\label{fig:cost_projections}
\end{figure}

\subsection*{Key Deployment Characteristics}

The key deployment characteristics are:

\begin{itemize}[leftmargin=0.5in]
  \item \textbf{RDS PostgreSQL instead of Aurora Serverless.} Standard RDS PostgreSQL \texttt{db.t3.micro} (PostgreSQL 16.6, 20GB gp2) replaces the originally proposed Aurora Serverless to stay within budget. The pgvector extension and HNSW indexing remain fully functional.
  \item \textbf{Groq as primary LLM backend.} Groq free-tier API (Llama 3.1 8B Instant) serves as the primary backend, with Amazon Bedrock (Nova Micro), Ollama, and HuggingFace as fallbacks. The auto-fallback chain ensures generation availability.
  \item \textbf{Default VPC with security groups.} The SAM template uses the default VPC with two subnets. RDS is configured as publicly accessible with security group restrictions; custom VPC migration is identified as a future improvement.
  \item \textbf{Fine-grained IAM roles.} Custom IAM role definitions provide least-privilege access for each Lambda function.
  \item \textbf{Cost governance.} Cost governance is implemented primarily through programmatic AWS Cost Explorer checks in the inference path, supplemented by CloudWatch metrics, SNS-based alerting, and backend fallback logic. Native AWS Budgets was considered but is not part of the active Free Tier SAM template; budget enforcement relies on the runtime Cost Explorer guard and the automatic LLM backend fallback chain.
  \item \textbf{LLM backend set to \texttt{groq}.} The \texttt{LLM\_BACKEND} environment variable defaults to \texttt{groq}, with Amazon Bedrock and Ollama retained as fallbacks.
\end{itemize}

All infrastructure is defined in a SAM template (\texttt{template-freetier.yaml}) and configuration file (\texttt{samconfig-freetier.toml}), ensuring reproducible and consistent deployments from the same codebase. Step-by-step reproduction instructions are provided in Appendix~\ref{app:setup}.

\section{Test Coverage}
\label{sec:tests}

The system includes 208 unit tests distributed across 15 test modules covering all major subsystems: ingestion (5 sources), preprocessing (17 normalization steps), embedding, retrieval (BM25 caching, RRF fusion, age-days computation, cross-encoder reranking), RAG orchestration, LLM backend fallback, observability (CloudWatch metrics, drift detection), pipeline orchestration, queue management, storage, API endpoints (including /dashboard/insights smoke tests), scheduler, and data sync/migration utilities (\texttt{test\_sync\_to\_aws.py}). All tests execute without external dependencies (database, AWS services) using mocked interfaces where appropriate.

\section{Experiment Tracking and Result Versioning}
\label{sec:experiment_tracking}

Experiment results are tracked through structured JSON artifacts versioned alongside the source code in Git. Each evaluation run produces three output files: \texttt{eval\_summary.json} (aggregate metrics, latency statistics, and cost projections for both configurations), \texttt{ragas\_scores.json} (per-query RAGAS faithfulness, answer relevancy, and context precision scores for all 100 observations), and \texttt{raw\_results.json} (complete per-query records including retrieved contexts, generated answers, citation grounding details, and timing breakdowns). This file-based approach was chosen over a dedicated experiment tracking platform (e.g., MLflow, Weights \& Biases) for three reasons: (1) the evaluation compares exactly two fixed configurations (baseline vs.\ hybrid) rather than a large hyperparameter sweep requiring run management; (2) JSON artifacts are directly consumable by the statistical analysis and visualization scripts without additional API dependencies; and (3) Git versioning provides full traceability of which code revision produced which results, with \texttt{git diff} enabling precise identification of changes between evaluation runs. The design-phase grid search results (\texttt{evaluation/stale-results-hybrid/}) are similarly retained under version control to document the design evolution from linear scoring to RRF-based fusion.

\section{Custom CloudWatch Metrics}
\label{sec:metrics}

Eight custom CloudWatch metrics have been implemented within a single \texttt{TechPulse/\{stage\}} namespace to provide operational visibility:

\begin{enumerate}[leftmargin=0.5in]
  \item \texttt{IngestionCount} --- documents ingested per source per cycle
  \item \texttt{PipelineChunks} --- chunks generated per pipeline run
  \item \texttt{PipelineLatency} --- end-to-end pipeline execution time
  \item \texttt{PipelineErrors} --- pipeline failure count
  \item \texttt{ApiLatency} --- per-request API response time
  \item \texttt{HallucinationFlags} --- count of citation grounding failures
  \item \texttt{DriftMeanSimilarity} --- mean cosine similarity from the 20-query probe set
  \item \texttt{DriftAlert} --- binary flag (0/1) indicating whether drift criteria are breached
\end{enumerate}

\section{Engineering Observations}
\label{sec:eng_obs}

Several non-obvious findings arose during implementation that informed the final design:

\begin{itemize}[leftmargin=0.5in]
  \item \textbf{Lambda cold-start and MiniLM.} Initial cold-start delays of 8--12\,s were resolved by a lazy-initialized singleton and 2{,}048\,MB memory allocation for the RAG API function, reducing model load to $<$2\,s on warm invocations. A 5-minute EventBridge warmup ping (\texttt{ScheduleV2}) keeps the RAG Lambda warm by sending a lightweight \texttt{\{"httpMethod":"WARMUP"\}} event that the handler detects and returns immediately, avoiding cold-start latency for user-facing queries.

  \item \textbf{SHA-256 deduplication.} The \texttt{ON CONFLICT DO NOTHING} strategy proved fully effective; ArXiv produced the most hash collisions ($\sim$15--20\% per cycle) due to preprint updates.

  \item \textbf{LocalStack fidelity.} S3 and SQS operations matched production behavior, enabling integration testing without AWS credentials. CloudWatch metrics required cloud deployment to verify.

  \item \textbf{HuggingFace rate limits.} The free-tier rate limit ($\sim$30 req/min) triggered the Ollama fallback during sustained evaluation runs, confirming the fallback chain operates correctly.

  \item \textbf{pgvector HNSW index build.} Index rebuilds took 45--90\,s on db.t3.micro for corpora above $\sim$10k chunks---acceptable for batch but would need management for real-time updates.
\end{itemize}

\noindent These implementation decisions reflect practical trade-offs between performance, cost, and system complexity in a constrained cloud environment.

\section{Hybrid Retrieval: BM25 Index Caching \& Architectural Details}
\label{sec:bm25_caching}

The BM25 component of hybrid retrieval employs module-level index caching with automatic periodic refresh. The full-corpus BM25 index is constructed from all indexed chunks during the initial load and then cached in process memory. The cache is automatically invalidated and reconstructed every 30 minutes to incorporate newly ingested documents while minimizing per-query recomputation overhead. This caching strategy reduces mean BM25 retrieval latency by $\sim$200\,ms (per-query computation ~$O(N)$ tokens against the corpus) but introduces a bounded staleness: queries execute against corpus state from at most 30 minutes prior. For drift detection validation runs, probe queries execute against potentially cached BM25 results; sensitivity analysis confirms that cache age induces <1\% variation in mean cosine similarity measurements, well within the 10\% drift threshold margin.

The BM25 tokenizer (class \texttt{BM25Okapi}) is independent of the query processing tokenizer to ensure consistent vocabulary across corpus indexing and retrieval. Token filtering removes $\sim$110 domain-adapted stopwords (including common tech-domain terms like ``tech'', ``technology'', ``article'', ``data''), enforces minimum token length of 2 characters, and normalizes the text to lowercase before tokenization. This design choice prevents false positive matches on ubiquitous terms while preserving domain-specific signal.

\section{Ingestion Scheduler \& Orchestration}
\label{sec:scheduler}

TechPulse implements two ingestion orchestration modes: (1) AWS EventBridge-based scheduling for cloud deployments, and (2) local polling-based scheduler for Docker Compose environments.

\noindent\textbf{AWS EventBridge:} On AWS, a CloudWatch Events rule (cron expression \texttt{cron(0 */6 * * ? *)}) triggers the ingestion Lambda every 6 hours (00:00, 06:00, 12:00, 18:00 UTC). Each invocation sequentially calls all five ingesters (ArXiv, Hacker News, DEV.to, GitHub Trending, RSS) within a single Lambda function, minimizing cold-start overhead. The 6-hour cycle was determined empirically to balance corpus freshness (query latency impact $<$100\,ms for 6-hour-old documents in the tech domain) against AWS Lambda free-tier invocation limits (1 million/month, ~4 invocations/day).

\noindent\textbf{Local Scheduler:} In Docker Compose development, a dedicated \texttt{scheduler} service implements polling-based orchestration with source-specific refresh intervals: ArXiv every 12 hours (lower volume), and Hacker News, DEV.to, GitHub, RSS every 30 minutes (higher velocity signals). This asymmetric scheduling reflects domain knowledge about source update patterns.

\section{API Endpoints \& Request Handling}
\label{sec:api_endpoints}

The FastAPI application exposes four primary endpoints:

\begin{enumerate}[leftmargin=0.5in]
  \item \texttt{GET /health} --- Deep health check returning the status of database, S3, SQS, and configured LLM backend (response time: $<$2\,s). Returns HTTP~200 when all probes pass, or HTTP~503 with per-component status when any probe fails, enabling load-balancer and uptime-monitor integration.
  
  \item \texttt{POST /ask} --- Main query endpoint. Accepts \texttt{query} (string, 1--2000 characters), \texttt{mode} (``baseline'' or ``hybrid''), and optional \texttt{sources} (list of up to 5 source types to filter, e.g., \texttt{["arxiv", "hn"]}). Source values are validated against \texttt{VALID\_SOURCES = \{arxiv, hn, github, devto, rss\}}; invalid values return HTTP~422 with the allowed set. Returns ranked passages, generated answer (if LLM available), and citation metadata. The top-$k$ parameter is configured server-side via the \texttt{TOP\_K} environment variable (default~5) and is not exposed as a client parameter to prevent abuse. Enforced rate limit: 10 requests/minute per IP address (SlowAPI).
  
  \item \texttt{POST /drift} --- Manual drift detection trigger (rate-limited to 2 requests/minute). Runs the canonical 20-query probe set through hybrid retrieval and evaluates two criteria: relative quality drop $> 10\%$ and Shewhart $3\sigma$ lower control limit breach; returns mean/std similarity, baseline reference, and alert reason.
  
  \item \texttt{GET /dashboard/insights} --- Complex aggregation endpoint returning dashboard data: trending topics (keywords + growth metrics), 30-day mention time-series, source distribution, and live highlight snippets (top 5 by composite score). Optional query parameters: \texttt{days} (7--60, default 30) and \texttt{sources} (comma-separated filter over ARXIV, HN, DEVTO, GITHUB, RSS).
\end{enumerate}

Rate limiting is implemented via SlowAPI (dependency: \texttt{slowapi}) to provide basic DOS protection and fair resource allocation in shared environments. Per-endpoint limits are set according to computational cost: \texttt{/ask} at 10~req/min (each invocation incurs retrieval, optional cross-encoder reranking, and LLM generation) and \texttt{/drift} at 2~req/min (each invocation evaluates 20 probe queries against the drift detection criteria). Monitoring during evaluation runs showed peak throughput of $\sim$2--3 queries/min from a single user.

\section{Citation Grounding Verification}
\label{sec:citation_grounding}

The three-layer hallucination verification strategy includes an output-layer citation grounding check that validates the generated response's reference completeness. The verification logic:

\begin{enumerate}[leftmargin=0.5in]
  \item \textbf{Sentence splitting:} The response is split on sentence-ending punctuation using regex \texttt{r'(?<=[.!?]){\textbackslash}s+'}, producing atomic sentential units.
  
  \item \textbf{Citation detection:} Each sentence is checked for presence of citation markers matching the pattern \texttt{r'{\textbackslash}[Source{\textbackslash}s+{\textbackslash}d+{\textbackslash}]'}, where the captured index $N$ must satisfy $1 \leq N \leq \text{num\_sources}$.
  
  \item \textbf{Citation ratio computation:} The fraction $R = \text{cited\_sentences} / \text{total\_sentences}$ is computed and compared against a configurable threshold \texttt{CITATION\_GROUNDING\_THRESHOLD} (default 0.0 for local dev; deployed value 0.25 for production with Groq Llama~3.1 8B Instant, balancing citation enforcement against small-model compliance).
  
  \item \textbf{Hallucinated reference detection:} Invalid source indices (e.g., \texttt{[Source 12]} when only 5 sources are available) trigger a hallucination flag.
\end{enumerate}

Responses flagged as under-grounded (citation ratio $< $ threshold) or containing invalid references are recorded in CloudWatch as \texttt{HallucinationFlags} metrics for operational monitoring. During evaluation, the RAGAS-computed ``Citation Grounding'' metric (Section~\ref{sec:eval_metrics}) correlates with but is independent of this output-layer check, providing complementary signal for model quality assessment.

\section{System Demonstration}
\label{sec:demo}

This section presents screenshots from the deployed system, demonstrating end-to-end functionality across key components. A full screen-recording walkthrough of the live system is available at: \url{https://youtu.be/iQ6LFZnk8F8}.

\subsection{Frontend: Overview and Search Interface}

Figure~\ref{fig:demo_overview_search} presents the two primary frontend views. The landing page (left) features the system branding, call-to-action buttons, and a hero section introducing TechPulse. The search interface (right) provides the full-featured query panel with retrieval mode selection (Baseline or Hybrid), recent query history in the sidebar, and suggested starter queries.

\begin{figure}[htbp]
\centering
\begin{subfigure}[t]{0.48\textwidth}
\centering
\includegraphics[width=\textwidth]{techpulse-overviewpage.png}
\caption{Landing page.}
\label{fig:demo_overview}
\end{subfigure}
\hfill
\begin{subfigure}[t]{0.48\textwidth}
\centering
\includegraphics[width=\textwidth]{rag_overview.png}
\caption{Search interface with sidebar and query history.}
\label{fig:demo_search}
\end{subfigure}
\caption{TechPulse frontend: landing page (left) and main search interface (right).}
\label{fig:demo_overview_search}
\end{figure}

\subsection{Frontend: Hybrid vs.\ Baseline Retrieval Comparison}

Figure~\ref{fig:demo_compare} shows the updated frontend executing queries in both retrieval modes. The hybrid response (left) includes diverse, multi-source citations fused via RRF with relevance scoring, while the baseline response (right) relies solely on vector similarity. This figure is intended as a qualitative interface illustration; quantitative retrieval conclusions are based on the statistical results reported in Chapter~5.

\begin{figure}[htbp]
\centering
\begin{subfigure}[t]{0.48\textwidth}
\centering
\includegraphics[width=\textwidth]{hybridmodel.png}
\caption{Hybrid retrieval (BM25 + vector + RRF).}
\label{fig:demo_hybrid}
\end{subfigure}
\hfill
\begin{subfigure}[t]{0.48\textwidth}
\centering

\includegraphics[width=\textwidth]{baseline.png}
\caption{Baseline retrieval (vector-only).}
\label{fig:demo_baseline}
\end{subfigure}
\caption{TechPulse frontend: hybrid vs.\ baseline retrieval for the same query.}
\label{fig:demo_compare}
\end{figure}

Figure~\ref{fig:demo_compare_modes} shows the built-in compare view, which presents baseline and hybrid responses side-by-side for the same query, enabling direct qualitative comparison of answer quality, source diversity, and citation patterns.

\begin{figure}[htbp]
\centering
\includegraphics[width=0.95\textwidth]{compare model.png}
\caption{TechPulse frontend: side-by-side compare mode (baseline vs.\ hybrid) for the same query.}
\label{fig:demo_compare_modes}
\end{figure}

Figure~\ref{fig:demo_loading} illustrates the step-by-step loading sequence visible during hybrid retrieval, providing real-time progress feedback to the user as the system fetches documents, ranks evidence, and generates the grounded answer.

\begin{figure}[htbp]
\centering
\includegraphics[width=0.45\textwidth]{stepbysteploading_rag.png}
\caption{TechPulse frontend: step-by-step RAG pipeline loading indicator during hybrid retrieval.}
\label{fig:demo_loading}
\end{figure}

\subsection{Frontend: Dashboard \& System Monitoring}

In addition to the primary search interface, a comprehensive dashboard provides real-time system monitoring and evaluation results visualization. The dashboard aggregates the following functional components:

\begin{enumerate}[leftmargin=0.5in]
  \item \textbf{Trending Topic Analysis.} Displays the top-ranked keywords aggregated across all five data sources over a configurable time window (default: 30 days). Keyword extraction uses a 68-entry canonical phrase map (\texttt{PHRASE\_MAP}) that normalises variant spellings and compound terms to a single canonical form (e.g., ``large language model'' $\to$ ``LLM'', ``retrieval augmented generation'' $\to$ ``RAG''). Domain-specific stopwords ($\sim$110 terms) are excluded before frequency counting. For each topic, the interface reports weekly and monthly mention counts alongside source coverage metrics and growth percentage (relative change of daily mentions in the past~7 days versus the preceding period). Topics appearing in three or more distinct sources are flagged as \emph{cross-source buzz} signals.

  \item \textbf{Research-vs-Practice Gap Detection.} A supplementary analytic layer classifies each trending keyword by its dominant source profile: \textit{Research-heavy} topics appear predominantly in ArXiv; \textit{Practice-heavy} topics appear predominantly in Hacker News, DEV.to, and GitHub. This gap metric highlights emerging technologies that have strong academic attention but limited practitioner adoption, or vice versa.

  \item \textbf{Insight Template System.} To produce human-readable narrative summaries, trending keywords are mapped to one of seven trend categories (\texttt{rapid\_rise}, \texttt{steady\_rise}, \texttt{slight\_decline}, \texttt{steep\_decline}, \texttt{stable}, \texttt{research\_heavy}, \texttt{practice\_heavy}) with pre-authored natural-language templates. Template selection is deterministic per keyword via hash-based indexing (\texttt{hash(keyword) \% len(pool)}), ensuring consistent presentation across page refreshes.

  \item \textbf{Time-Series Trend Visualization.} A line chart tracking mention frequency for trending topics over 30 days, enabling stakeholders to observe temporal patterns in technology adoption.

  \item \textbf{Source Distribution Analytics.} A bar chart and tabular view showing document counts per ingestion source (ArXiv, Hacker News, DEV.to, GitHub, RSS), allowing corpus composition assessment.

  \item \textbf{Live Insight Explorer.} An interactive panel enabling source-based filtering of daily highlights (top 5 items by composite score per ingestion cycle). Composite scoring combines keyword overlap with current trending terms, content length, and a source-type boost (ArXiv documents receive $+0.4$ to prioritise peer-reviewed research). Highlights are presented as clickable links that resolve to source-specific search interfaces: ArXiv API search, GitHub repository search, Hacker News Algolia index, DEV.to tagged search, or Bing News feed search. Resolution logic (\texttt{\_resolve\_highlight\_url()}) falls back to topic-aware search queries when direct URLs are unavailable.

  \item \textbf{Pipeline Architecture Diagram.} A sequential six-stage flowchart visualizing the ingestion-to-generation pipeline (Data Sources → Ingestion → Embedding → pgvector DB → Hybrid Retrieval → LLM Answer Generation), providing non-technical stakeholders with system architecture intuition.

  \item \textbf{Evaluation Results Aggregation.} Displays all evaluation metrics computed in Chapter~5:
  \begin{itemize}
    \item Radar plot comparing baseline and hybrid RAGAS scores (five dimensions: Faithfulness, Answer Relevancy, Context Precision, Citation Grounding, Composite).
    \item Bar chart comparing mean and 95th percentile response latencies for both retrieval modes.
    \item Line chart of design-phase sensitivity analysis (superseded by wRRF), showing composite score variation across $\alpha \in [0.0, 1.0]$ (semantic weight) under the linear-reranking scheme that was replaced by weighted RRF in the final system.
    \item Summary statistics cards: sample count (100), hybrid composite score (0.723), composite gain (+4.7 pp), and projected monthly infrastructure cost (\$0 infrastructure at Free Tier; total imputed cost including Groq LLM inference is $\sim$\$0.33/month at 100~queries/day, as reported in Table~\ref{tab:cost_breakdown}).
  \end{itemize}

  \item \textbf{Multi-Source Attribution.} All highlights include explicit source attribution and clickable source chips at the dashboard base, allowing rapid filtering by data origin and validation that the system ingests from all five specified sources.
\end{enumerate}

The dashboard serves two purposes: (1) communicating system capabilities and results to non-technical stakeholders (e.g., project advisors, peer reviewers) through summary statistics and visualizations, and (2) enabling real-time operational monitoring of ingestion freshness and source distribution balance. Interactive features (source filtering, clickable external links, animated metrics) provide immediate feedback to users exploring the system's outputs.

Figures~\ref{fig:demo_dashboard_topics}--\ref{fig:demo_dashboard_eval} present the deployed dashboard interface.

\begin{figure}[htbp]
\centering
\begin{subfigure}[t]{0.48\textwidth}
\centering
\includegraphics[width=\textwidth]{techtopichighlight.png}
\caption{Tech Topic Highlights: trending topics with growth metrics, source coverage, and weekly/monthly mention counts.}
\label{fig:demo_dashboard_topics}
\end{subfigure}
\hfill
\begin{subfigure}[t]{0.48\textwidth}
\centering
\includegraphics[width=\textwidth]{analytics_liveinsight.png}
\caption{Analytics \& Live Insight Explorer: 30-day mention trend, source distribution, and daily highlights with source filtering.}
\label{fig:demo_dashboard_insights}
\end{subfigure}
\caption{TechPulse dashboard: trending topic analysis (left) and analytics with live insight explorer (right).}
\label{fig:demo_dashboard_main}
\end{figure}

\begin{figure}[htbp]
\centering
\begin{subfigure}[t]{0.48\textwidth}
\centering
\includegraphics[width=\textwidth]{evalresult_1.png}
\caption{RAGAS metrics radar (Baseline vs.\ Hybrid) and response latency comparison (mean and P95).}
\label{fig:demo_dashboard_eval1}
\end{subfigure}
\hfill
\begin{subfigure}[t]{0.48\textwidth}
\centering
\includegraphics[width=\textwidth]{evalresult_2.png}
\caption{Hybrid weight sensitivity analysis and summary statistics (sample count, composite score, gain, monthly cost).}
\label{fig:demo_dashboard_eval2}
\end{subfigure}
\caption{TechPulse dashboard: evaluation results --- RAGAS metrics and latency (left), sensitivity analysis and key statistics (right).}
\label{fig:demo_dashboard_eval}
\end{figure}

\subsection{Deployment Validation: CI/CD, Infrastructure, and Database}

Three key deployment components have been validated and are operational:

\noindent\textbf{(1) CI/CD Pipeline (GitHub Actions):} Figure~\ref{fig:demo_cicd} (left) shows the GitHub Actions workflow executing successfully with all lint, test, build, and deployment stages passing.

\noindent\textbf{(2) Infrastructure (AWS CloudFormation):} Figure~\ref{fig:demo_cfn} (center) shows the deployed AWS CloudFormation stack with all resources provisioned: Lambda functions, API Gateway, S3 buckets, SQS queues, EventBridge rules, and IAM roles.

\noindent\textbf{(3) Database (RDS PostgreSQL):} Figure~\ref{fig:demo_rds} (right) shows the AWS RDS console confirming the TechPulse PostgreSQL instance is running with pgvector extension installed.

\begin{figure}[htbp]
\centering
\begin{subfigure}[t]{0.30\textwidth}
\centering
\includegraphics[width=\textwidth]{github workflow run.png}
\caption{CI/CD: GitHub Actions workflow.}
\label{fig:demo_cicd}
\end{subfigure}
\hfill
\begin{subfigure}[t]{0.32\textwidth}
\centering
\includegraphics[width=\textwidth]{techpulse-cloudformation.png}
\caption{Infrastructure: CloudFormation stack.}
\label{fig:demo_cfn}
\end{subfigure}
\hfill
\begin{subfigure}[t]{0.32\textwidth}
\centering
\includegraphics[width=\textwidth]{rds-techpulse.png}
\caption{Database: RDS PostgreSQL instance.}
\label{fig:demo_rds}
\end{subfigure}
\caption{Deployment validation: CI/CD pipeline, CloudFormation stack, and RDS database all operational.}
\label{fig:demo_deployment}
\end{figure}

\medskip
\noindent This chapter documented implementation status across 20 components and key engineering observations. All core pipeline components are integration-tested; cloud deployment via SAM is validated. Chapter~7 presents conclusions, limitations, and future work.

